\section{问题三：两种代表工况策略与调节优先级}

\subsection{代表工况选择}

以入口质量负荷 $C_{\mathrm{in}}Q$ 为主轴，在样本量充足的稳定状态中选择 R2 与 R8。R2 的平均入口浓度为 $25.843\gNm$、质量负荷为 $1.232\times10^7$；R8 的平均入口浓度为 $46.090\gNm$、质量负荷为 $2.217\times10^7$，后者高约 79.9\%。二者中位驻留时间分别为 852 min 与 870 min，排除了极短异常段造成的偶然差异。

\subsection{策略参数表}

表\ref{tab:q3u} 和表\ref{tab:q3t} 将最优参数与该工况历史均值并列。其目的不是要求现场一步跳到目标值，而是给出工况稳定后可采用的连续设定中心；实际控制应增加斜率限制和切换滞回，并按 $0.1\kV$、$1\mathrm s$ 的设备分辨率取整后重新复核。

\begin{table}[H]
\centering
\caption{低、高负荷代表工况的电压策略}
\label{tab:q3u}
\resizebox{\textwidth}{!}{%
\begin{tabular}{ccrrrrrrrr}
\toprule
工况 & 标签 & $U_1^*$ & 历史 $U_1$ & $U_2^*$ & 历史 $U_2$ & $U_3^*$ & 历史 $U_3$ & $U_4^*$ & 历史 $U_4$\\
\midrule
低负荷 & R2 & 63.107 & 60.496 & 68.244 & 60.522 & 57.538 & 51.338 & 58.165 & 51.323\\
高负荷 & R8 & 70.430 & 62.055 & 69.267 & 62.125 & 58.431 & 44.572 & 57.505 & 44.633\\
\bottomrule
\end{tabular}}
\end{table}

\begin{table}[H]
\centering
\caption{低、高负荷代表工况的振打与风险指标}
\label{tab:q3t}
\resizebox{\textwidth}{!}{%
\begin{tabular}{ccrrrrrrrrrr}
\toprule
工况 & $T_1^*$ & 历史 $T_1$ & $T_2^*$ & 历史 $T_2$ & $T_3^*$ & 历史 $T_3$ & $T_4^*$ & 历史 $T_4$ & $P^*$ & $q_{.95}(C)$\\
\midrule
R2 & 290.0 & 247.7 & 289.0 & 247.6 & 506.0 & 443.9 & 510.0 & 443.5 & 1961.4 & 9.503\\
R8 & 249.6 & 195.9 & 249.7 & 195.4 & 440.8 & 440.6 & 435.4 & 440.4 & 2151.9 & 9.226\\
\bottomrule
\end{tabular}}
\end{table}

低负荷 R2 的尘负荷指数仅 1.098，四场周期均可延长至历史上界以降低振打频率；高负荷 R8 的指数达到 1.45，前两场周期约为 250 s，第三、第四场分别约为 441 s 与 435 s，体现尘负荷活跃时的周期收紧。电压方面，两类工况均采用四场协同分配；高负荷 R8 尤其需要提高末两场相对历史水平，不能把捕集任务固定压在某一单场。

\subsection{边际收益与优先级切换}

图\ref{fig:priority} 展示式\eqref{eq:priority} 的反事实结果。低负荷 R2 以 $U_1,U_3$ 居前；高负荷 R8 的四场收益十分接近，当前点以 $U_4,U_2$ 略高。电压项整体显著高于周期项，但这不意味着振打不重要：周期主要承担尘负荷和瞬时峰值安全约束，其价值不能只由稳态浓度斜率衡量。

\begin{figure}[H]
\centering
\includegraphics[width=0.74\textwidth]{fig10_priority_comparison.png}
\caption{低、高负荷工况的单位新增功率减排收益}
\label{fig:priority}
\end{figure}

低负荷 R2 中，电压收益从高到低为 $U_1$（0.0614）、$U_3$（0.0610）、$U_2$（0.0565）、$U_4$（0.0559），单位为 $\mathrm{mg/(Nm^3\!\cdot\!kW)}$；最高的周期收益 $T_2$ 为 0.0114。高负荷 R8 中，$U_4,U_2,U_1,U_3$ 分别为 0.0586、0.0572、0.0570、0.0562，差距不足 4.3\%，说明严格工况下应按实时边际收益协同分配，而不能预设固定场序；周期项直接收益较小，但尘负荷约束已经活跃。

据此形成表\ref{tab:rules} 的优先级规则。

\begin{table}[H]
\centering
\caption{工况—风险—优先动作控制规则}
\label{tab:rules}
\begin{tabularx}{0.96\textwidth}{C{3.0cm}C{3.2cm}L L}
\toprule
识别信号 & 主导风险 & 第一动作 & 第二动作\\
\midrule
低/中负荷，$I_M<1.30$ & 基线穿透 & 按 $E_j$ 提高高收益电场电压 & 在风险裕量内延长周期节能\\
高负荷，$I_M\rightarrow1.45$ & 积灰与峰值 & 缩短活跃电场周期并错峰 & 按 $E_j$ 协同分配四场电压，保留快速补偿余量\\
高温异常 & 迁移能力下降、反电晕不确定 & 前后场重分配，检查火花率 & 避免仅靠单场堆高电压\\
工况切换边界 & 频繁切换与过冲 & 入口质量负荷前馈+双阈值滞回 & 电压斜率限制与周期缓变\\
短时超限 & 多场峰值叠加 & 立即禁止同步振打 & 临时提高末级电压并进入应急模式\\
\bottomrule
\end{tabularx}
\end{table}

\subsection{现场切换逻辑}

建议以 $15$ 分钟稳健工况标签作为慢层，以分钟级入口质量负荷作为快层。进入高负荷状态的阈值应高于退出阈值，避免在聚类边界反复跳变；电压目标通过斜率限制逐步到达，振打周期只在一个完整振打循环结束后更新。若某场尘负荷风险高于阈值，则周期安全优先级高于电压经济性优先级。四场的振打触发相位错开，末级电场保留快速补偿通道。
